\documentclass{ctexart}
\usepackage{graphicx} % Required for inserting images
\usepackage[margin=2.5cm]{geometry}
\usepackage{booktabs}
\usepackage{float}
\usepackage{soulutf8}
\usepackage{enumitem}
\usepackage{pifont}
\usepackage{pgfplots}
\usepackage{longtable}
\pgfplotsset{compat=newest}
\usepackage{amssymb,amsthm}% 数学写作环境：定义/定理/引理/例题/解答，共用计数器并按节编号
\renewcommand{\proofname}{\normalfont\bfseries 证明}
\newtheoremstyle{cnthm}{0.8em}{0.8em}{\normalfont}{}{\bfseries}{.}{0.5em}{}
\theoremstyle{cnthm}
\newtheorem{theorem}{定理}[section]
\newtheorem{lemma}[theorem]{引理}
\newtheorem{definition}[theorem]{定义}
\newtheorem{example}[theorem]{例题}
\newtheorem{solution}[theorem]{解答}
\usepackage{tikz}
\usepackage{xcolor} % 用于颜色设置
\usepackage{listings} % 用于代码块
\usepackage{fontspec} % 关键：用于使用系统字体 (XeLaTeX/LuaLaTeX)
% Linux 环境中没有 Windows 宋体时，使用可用的中文衬线字体作为宋体替代
\setCJKmainfont{Noto Serif CJK SC}
\usepackage{authblk}
\usepackage[
    colorlinks=true,      % 使用颜色而非方框
    linkcolor=blue,       % 内部链接颜色
    citecolor=green,      % 引用链接颜色
    urlcolor=red,         % URL 链接颜色
    bookmarks=true,       % 生成 PDF 书签
    bookmarksopen=true,   % 默认展开书签
    pdfstartview=FitH     % 页面宽度适应窗口
]{hyperref}

\sethlcolor{yellow}%设置所需高亮颜色
\usepackage[most]{tcolorbox}
\tcbuselibrary{listings,skins,breakable}
% 页眉页脚
\usepackage{fancyhdr}
\pagestyle{fancy}

% 定义草绿色
\definecolor{grassgreen}{RGB}{124,179,66}
% 补全缺失的浅蓝色定义（防止编译报错）
\definecolor{lightblue}{RGB}{30,144,255} 

% 清空默认
\fancyhf{}
\setlength{\headheight}{14pt}

% 页眉
\fancyhead[L]{\textcolor{lightblue}{\kaishu Computer Vision}}
\fancyhead[R]{\textcolor{lightblue}{\textit{2026.9-2027.1}}}

% 页脚页码居中
\fancyfoot[C]{\textcolor{lightblue}{\thepage}}

% 页眉、页脚线颜色
\renewcommand{\headrulewidth}{1pt}
\renewcommand{\footrulewidth}{1pt}
\renewcommand{\headrule}{\color{blue!30!white}\hrule height 1pt width\headwidth}
\renewcommand{\footrule}{\color{blue!30!white}\hrule height 1pt width\headwidth}
\newtcolorbox{fancybox}[2][blue]{ % 可选颜色参数
  enhanced,
  breakable,
    colback=#1!3!white,
    colframe=#1!18!white,
    boxrule=0.8pt,
    arc=6pt,
    outer arc=6pt,
  left=10pt,right=10pt,top=12pt,bottom=10pt,
  title={\textbf{#2}},
    fonttitle=\bfseries\color{#1!65!black},
  attach boxed title to top left={yshift=-2mm,xshift=5mm},
  boxed title style={
        colback=#1!8!white,
        colframe=#1!18!white,
        boxrule=0.6pt,
        arc=4pt,
  },
    shadow={1.5pt}{-1.5pt}{0pt}{black!10},
}

% 定义一套适合C++的配色方案
\usepackage{tabularx}
\newfontfamily{\consolasfont}{Ubuntu}

\newcommand{\bluecode}[1]{{\color{blue}\consolasfont #1}}
\lstdefinestyle{MyCppStyle}{
    language = C++,
    % 确保这里使用 \consolasfont
    basicstyle = \consolasfont\small, 
    keywordstyle = \color{blue!70!black}, 
    commentstyle = \color{green!50!black}, 
    stringstyle = \color{red!60!black}, 
    numberstyle = \tiny\color{gray}, 
    numbers = left, 
    frame = single, 
    rulecolor = \color{lightgray}, 
    breaklines = true, 
    backgroundcolor = \color{white}, 
    tabsize = 4, 
    showstringspaces = false, 
    xleftmargin = 2em, 
    escapeinside=``,
}
% --- lstdefinestyle HTML 风格定义 ---
\lstdefinestyle{mhtml}{
    language = HTML,
    basicstyle = \consolasfont\small, 
    tagstyle = \color{blue!70!black}\bfseries,
    keywordstyle = \color{purple!80!black},
    commentstyle = \color{green!50!black}\itshape, 
    stringstyle = \color{red!60!black}, 
    numberstyle = \tiny\color{gray}, 
    numbers = left, 
    frame = single, 
    rulecolor = \color{lightgray}, 
    breaklines = true, 
    backgroundcolor = \color{white}, 
    tabsize = 4, 
    showstringspaces = false, 
    xleftmargin = 2em, 
    escapeinside=``,
    morekeywords={als, doctype, html, head, title, meta, link, body, div, span, h1, h2, h3, h4, h5, h6, p, a, img, ul, ol, li, table, tr, th, td, pre, code, script, style},
    morekeywords=[2]{class, id, href, src, rel, type, charset, lang, name, content, style},
    keywordstyle=[2]\color{purple!80!black},
}
\usepackage{hyperref}
\hypersetup{
    colorlinks=true,      
    linkcolor=blue,       
    filecolor=magenta,    
    urlcolor=cyan,        
    pdftitle={我的文档标题}, 
}

% ==================== 正文部分 ====================
\begin{document}

% 开启无页眉页脚模式（针对封面和目录）
\pagestyle{empty}

% ------------------ 美化封面页 ------------------
\begin{titlepage}
    \centering
    % 顶部 TikZ 几何几何装饰条（全宽 bleed 效果）
    \begin{tikzpicture}[remember picture, overlay]
        \fill[blue!50!white] (current page.north west) rectangle ([yshift=-2cm]current page.north east);
        \fill[grassgreen] (current page.north west) ++(0,-2cm) rectangle ([yshift=-2.2cm]current page.north east);
    \end{tikzpicture}
    
    \vspace*{3.5cm}
    
    % 大标题
    {\fontsize{30pt}{36pt}\selectfont\bfseries\color{blue!80!black} 计算机视觉（CV）}\\[0.6em]
    % 副标题或英文标识
    {\fontsize{14pt}{18pt}\selectfont\color{gray}\textsf{Computer Vision}}\\[1cm]
    
    % 居中装饰线
    \centering\makebox[4cm]{\color{grassgreen}\rule{5cm}{1.5pt}}
    
    \vfill
    
    % 使用 tcolorbox 包装的精致作者信息框
    \begin{tcolorbox}[
        enhanced,
        width=0.55\textwidth,
        colback=blue!5!white,
        colframe=blue!5!white,
        boxrule=1.2pt,
        arc=6pt,
        sharp corners=downhill, % 别致的对角切角效果
        shadow={2pt}{-2pt}{0pt}{black!20},
        halign=center,
        valign=center,
        top=15pt, bottom=15pt
    ]
        {\large\bfseries\color{black!80} Author：} {\large\kaishu Simson} \\[0.6em]
        {\large\bfseries\color{black!80} Date：} {\large\kaishu September 2026}
    \end{tcolorbox}
    
    \vfill
    \vspace*{2cm}
    
    % 底部标识
    {\color{gray!80}\small\the\year\ \copyright\ All Rights Reserved.}
\end{titlepage}

\pagestyle{fancy}

\newpage
\tableofcontents
\newpage

% 设置从当前页（正文第一页）开始采用阿拉伯数字编码，并自动重置为第 1 页
\pagenumbering{arabic}

\section{导论}

\textbf{可视计算的定义：}针对视觉内容生成与处理的所有计算机课学分制的总称（计算机图形学，计算机视觉，可视化，AR,VR）人类智能在感知，认知，决策，行动，感知的循环中逐渐积累。

重要课题：
\begin{itemize}
    \item 物理世界仿真与模拟
    \item 计算机图形学
    \item 计算机视觉
\end{itemize}

\section{VCL基础知识和技能}

\subsection{颜色，颜色感知和可视化}

人感知颜色靠的是对不同波长的光敏感的三种视锥细胞，这提供了颜色离散表达的基础。三种感光开展出RGB三色系统。颜色的匹配函数如下，是一个积分的办法。

\[
R=\int_0^{\infty}I(\lambda)\bar{r}(\lambda)d\lambda
\]
\[
G=\int_0^{\infty}I(\lambda)\bar{g}(\lambda)d\lambda
\]
\[
B=\int_0^{\infty}I(\lambda)\bar{b}(\lambda)d\lambda
\]

三种通道和原色的光由离散合成有两种机制：基于发光叠加相加混光（显示器发光原理）和基于白光吸收的相减混光（涂料和打印原理）。

计算机视觉基于人对光的感知。在人的感知之下，红色有前进的效应，蓝色会有看起来后退的效应。此即所谓色彩立体视觉。同时在人的认知下颜色也有冷色和暖色之分（这真的是课堂PPT里面的）。颜色信息在人眼中具有恒常性，不论在何种光照之下人对于颜色的认知不会发生大的改变。

当然这种逻辑感知也会带来相应的视觉错觉，譬如色诱导，使用补色可以让人自动脑补对应的颜色；再比如说边缘效应和移动错觉（动态立体视觉）。

人眼对于不同的光敏感性不同，对不同波段的光，施加灰度和模糊处理可以得到不同的认知效果（灰度指的是介于黑白之间的明暗程度）。

\section{显示（Display）}

\subsection{显示器的种类和发展历史}

\begin{enumerate}
    \item 矢量显示（阴极射线管）
    \item 二维显示（电子束）
    \begin{itemize}
        \item 像素：每个可以点亮的屏幕点
        \item 分辨率：品目像素格式
        \item 纵横比：像素列数除以行数
        \item 帧：显示的一个画面
        \item 隔行扫描：保证亮度均衡
    \end{itemize}
    \item 彩色二维显示（彩色荧光幕，不同荧光剂分区涂布）
    \item 二维显示（LCD）：统一发光，用液晶部分控制光的偏振，调节亮度
    \item 二维显示（OLED）：有机发光半导体显示器。独立的LED单元。
\end{enumerate}

\subsection{色域}

美国国家电视标准委员会的色域标准称为NTSC。

\begin{itemize}
    \item sRGB：微软，HP（72% NSTC）
    \item ARGB: Adobe颜色（设计软件）
    \item Rec.2020:4K(电视的标准)
    \item DCI-P3 : Apple标准
\end{itemize}

\subsection{像素标准}

显卡的存储和刷新越厉害，达到4K8K就越容易。

\subsection{立体显示}

立体显示依据的原理是双目视觉。所有的立体显示技术基于人类的生理特征。最开始的显示技术是双目偏振片的3D显示。沉浸式的3D显示技术（AR,VR,XR）跟踪头显，裸眼3D+眼睛跟踪。现有的技术向MR-mix方向进行发展。

裸眼3D技术：LCD屏幕的方法：微柱透镜显示技术。利用视差屏障，令OLED屏幕上不同的像素点用不同角度进入人眼。

全息显示技术（Hplographic display）：用干涉记录波前的信息（数字全息摄影），用衍射记录物光波场（受限于较小的范围）。

多模态显示技术：触觉显示（触觉输出），嗅觉显示（控制挥发的成分。。。？）

\section{图像（Image Representation）}

在二维空间中，每个点可以由四个参数表示。$E(x,y,\lambda,t)$其中$\lambda$字母的含义为光的波长。从数学上来说图像本身是一个连续的表征（真实世界），在屏幕上离散化显示。VCL的任务在于用离散方式进行还原。

\subsection{图像的表示}

图像显示的两个方法是矢量表示法（Vector Representation，对应SVG格式）和光栅化表示法（Raster Representation，像素级存储）。

矢量表示法的精确度高，用较小的空间，不会受到分辨率影响（可以用新的矢量计算适应不同屏幕不同大小）多用于图标等规则图形。Raster存储空间较大，受限于分辨率，但是可以表达更加不规则的图形。

\subsection{图形处理器（GPU）}

GPU有成百上千个简单的计算单元，可以在同一时间并行运行大量程序，比如计算每个像素的颜色，而CPU一般只有个位数的可并行的计算单元。故而GPU显卡适合处理进行大量计算的图形编程任务。

Pixel存储中有重要概念帧缓冲区（Frambuffers for image processing and display），指的是储存二维像素矩阵的内存大小。硬件以60Hz速度扫描帧缓冲区中的内容，这也决定了显示器的种类。

同样存储彩色图片的时候，Fullcolor（RGB），决定了显示器的颜色存储位数。一般来说真彩色图像的显示屏为24-bit位存储。如果颜色粗出空间较小，色彩的丰富度就会下降。

gif颜色存储的性质：Colormap映射。对于每一个像素点，存储一个8-bit常数，映射到RGB颜色表格Colormap，根据映射进行显示。在预先选好的表格中进行查找和压缩。

\begin{fancybox}{几种图片的文件格式}

\textbf{常见图像文件格式的全称：}
\begin{itemize}[leftmargin=1.6em,itemsep=1pt,topsep=3pt,parsep=0pt]
    \item \textbf{BMP}：Full color image（全彩色图像）
    \item \textbf{JPEG}：Joint Photographic Experts Group Format
    \item \textbf{TIFF}：Tagged-Image File Format
    \item \textbf{GIF}：CompuServe Graphics Interchange Format
    \item \textbf{PPM}：Portable PixMap Format（ASCII 或二进制）
    \item \textbf{EPS}：Encapsulated PostScript Format（ASCII）
\end{itemize}

\vspace{2pt}
\begin{center}
\small
\begin{tabular}{lccc}
\toprule
 & \textbf{BITS PER PIXEL} & \textbf{FILE SIZE} & \textbf{COMMENTS} \\
\midrule
JPEG & 24            & small  & lossy compression \\
TIFF & 8, 24         & medium & good general purpose \\
GIF  & 1, 4, 8       & medium & popular, but 8-bit \\
PPM  & 24            & big    & easy to read/write \\
EPS  & 1, 2, 4, 8, 24 & huge  & good for printing \\
\bottomrule
\end{tabular}
\end{center}
\end{fancybox}

同时，相较于24-bit的图像，有的高级图像会膨胀到每个像素96-bit，有一些其他的通道来表示图像。

\begin{enumerate}
    \item Alpha通道（8-bit）：透明度（Transparency），存储透明度常数，存在前中后背景。
    
    根据参数$\alpha$计算像素点的表示：$b = (1-\alpha) a_1 + \alpha a_2$
    \item Z-buffer（16-bit）：保存像素点所对应的3D空间点的深度。判断在叠加一层的时候应该占据多少通道和Alpha信息。
    
    \item 双通道表示：既要保存图像的现有数据，又要保存图像的旧信息，方便修改。
\end{enumerate}

\subsection{图像处理（Image Processing）}

图像处理按操作的作用范围可以分为两类基本操作。\textbf{点处理（Point Processing）}逐像素地修改像素值，输出仅是该像素自身像素值的函数；\textbf{滤波（Filtering）}的输出则是该像素周围（通常是局部）邻域（local neighborhood）内像素的函数，因而需要用到邻域中其他像素的信息。
我们也可以理解为一个是逐点操作，一个是区域化的操作（图像抖动，图像过滤，区域补全）。

除这两类基本操作之外，图像处理还包括以下其他课题：

\begin{itemize}
    \item 图像变换（Image transformation）：缩放（resize）、形变（warp）
    \item 图像压缩（Image compression）：无损（lossless）与有损（lossy）
    \item 图像编解码（Image encoding and decoding）：用于高效传输（efficient transmission）
    \item 图像分割（Image segmentation）
    \item 图像抠图（Image matting）
    \item 图像理解（Image understanding）：进一步走向计算机视觉（computer vision）
    \item 视频处理（Video processing）：推广到图像序列（a sequence of images）
    \item ……
\end{itemize}

\subsubsection{Dithering}

\paragraph{灰度图像频率的定义}

灰度图像频率可以理解为灰度数值变化的剧烈程度，数学表示通过对像素矩阵进行DFT（傅里叶变换）。

设 $f(x,y)$ 为 $M \times N$ 的灰度图像，其二维DFT及逆变换为：
\[
F(u,v) = \sum_{x=0}^{M-1}\sum_{y=0}^{N-1} f(x,y)\, e^{-j 2\pi \left(\frac{ux}{M} + \frac{vy}{N}\right)}
\]
\[
f(x,y) = \frac{1}{MN}\sum_{u=0}^{M-1}\sum_{v=0}^{N-1} F(u,v)\, e^{j 2\pi \left(\frac{ux}{M} + \frac{vy}{N}\right)}
\]

其中 $F(u,v)$ 为频域系数（$j$ 为虚数单位），$(u,v)$ 为频率坐标：$u,v$ 靠近原点对应低频，即灰度缓变区域；远离原点对应高频，即边缘和细节。变换本身不损失信息，逆变换可无损还原 $f(x,y)$。

\paragraph{Dithering和噪声种类}
抖动和量化研究的是有限空间下图片的表示，比如要把灰度图像表示成二值的图像。此时如果直接使用0.5为分界划分每个pixel。这时要给整体图片接一个整体的随机噪音（random dithering）。原来
是结构化的容易被人察觉的低频误差转换成频谱较为均匀的高频误差。（所谓频率是指的是低频区间的扰动）

施加噪音的方法有两种：Random Noise和Blue Noise，Blue Noise将误差的能量集中到高频的空间，低频范围内噪声较小，高频噪声能量较高，能够得到更好的表示效果。

噪声的施加如果作为输入，还会用于图像的模糊处理（Gaussian blur），均匀处理掉细节的误差之后Blue Noise会有更好的效果。

\paragraph{Dithering in Space}

用空间换表示效果：可以用多个像素点空间表示不同灰度，这时可以将抖动的幅度变成1-10的连续序列。所有的输入像素对应一个输出的pattern。这里
0-9的9像素显示模式是经过设计的，为了人尽量难以注意到细小的细节错误。这种方法叫作Ordered Dithering。

数学上的方法就是使用存储抖动矩阵记录抖动的序列，依次按照输入的归一化数值点亮像素点。

如果保持连续性不变化，不改变显示分辨率，进行有序抖动，可以用一个DIther Matrix将图片扫描一遍，像素点大于抖动阵的数值则保留，小于则纯白，可以实现同样分布的抖动。

\paragraph{Error Diffusion}

每一次抖动都会制造误差，可以通过这样的算法让每一次的误差逐步扩散到周边像素。Floyd算法可以决定抖动误差如何进行分配，让之前和之后的抖动误差有较好的平衡性和连续性。

如何在不同的建模和渲染场景中生成不同的合适的Blue noise是一个重要的话题，在几乎所有的图像处理场景中Dithering都是必要的。

\subsubsection{Filtering}

Filtering和图像的特征处理高度相关，最广泛和最简单的应用是Blurring模糊。让高频信息逐渐减少，剩下的集中在低频信息。
基本的Blurring Filter使用模糊核将图像的局部进行适当的归一化。

最简单的模糊核是一个3*3矩阵，无限维的卷积核相当于连续空间上的函数卷积

\[
h(t) = f * g = \int_{-\infty}^{\infty} f(\tau)g(t - \tau) d\tau
\]

相当于首先镜像然后向右平移。离散情况中卷积核在二维的像素图像上进行移动，常用的卷积核Laplace Kernel
在二维图像上卷积运算的数学表达如下，其中h为卷积核，b为输出，a为输入。

\[
b[x, y] = \sum_{u = -\infty}^{\infty}\sum_{v = -\infty}^{\infty}a[u,v]h[x - u,y - v]  
\]

Filtering方法可以进行图像特征提取，可以进行竖直方向和水平方向的有限差分，从而提取不同的信息。边提取算子的数学表达如下：

\[
\nabla a  = (\frac{\partial a}{\partial x}, \frac{\partial a}{\partial y}) , \qquad |\nabla a| = \sqrt{(\frac{\partial a}{\partial x})^2 + (\frac{\partial a}{\partial y}^2)}
\]

将以上计算梯度的梯度信息取整数近似可以得到Sobel edge Filter，是一个离散的非线性卷积核。
\[
\frac{\partial}{\partial x} \Rightarrow
\begin{bmatrix}
    -1 & 0 & 1 \\
    -2 & 0 & 2 \\
    -1 & 0 & 1
\end{bmatrix}
\]

有限差分可以解决一阶微分，同时可以用Laplace算子求二阶微分提取出图像的边缘。离散化的laplace算子意思是显然的，就是对比某个像素点和周边的差异。
laplace算子的连续近似如下：

\[
\nabla^2 f = \frac{\partial^2 f}{\partial x^2} + \frac{\partial^2 f}{\partial y^2}
\]

\[
\begin{aligned}
\frac{\partial^2 f}{\partial x^2} &\approx f(x+1, y) - 2f(x, y) + f(x-1, y) \\
\frac{\partial^2 f}{\partial y^2} &\approx f(x, y+1) - 2f(x, y) + f(x, y-1)
\end{aligned}
\]

对应的离散卷积核为：

\[
\mathbf{L} = \begin{bmatrix}
    0 & 1 & 0 \\
    1 & -4 & 1 \\
    0 & 1 & 0
\end{bmatrix}
\]

\subsubsection{Image inpainting}

图像补全的中心任务是填补image的某些缺失部分，让结果尽可能合理。这里的合理（Plausible）有以下含义：

\begin{itemize}
    \item Spatial locality ：和周围和谐
    \item Occam’s Razor ：不添加额外的非必要信息
    \item 梯度尽可能小
\end{itemize}

这种算法和计算过程常常称为laplace图像补全。这里优化函数的数学表示是

\[
f = \arg \min_f \iint ||\nabla f|| d\Omega
\]

这种技术还被应用于Image cloning/insertion，图像克隆和拼接，算法可以完成前景和背景的融合。相比于之前的Laplace处理，这里
要应用Poission处理，目标是将另一张图片注入另一张图片，优化函数令当前的梯度和图片尽可能接近。前面的Laplace处理是$\nabla g = 0$的特例
(背景是泛函分析中函数的变分法)

\[
Purpose: \min_f \iint_{\Omega} ||\nabla f - \nabla g|| ^ 2 dxdy
\]

在计算机科学中本质是离散意义上Possion和Laplace方程的求解。将梯度转换为差分，得到以下方程：

\[
\frac{f_{i+1,j} + f_{i-1,j} - 2f_{i,j}}{h^2} + \frac{f_{i,j+1} + f_{i,j-1} - 2f_{i,j}}{h^2} = \rho^2
\]

证明只需使用二元函数的Taylor展开，使用差分的方式进行计算，保留一次项。在实际应用中会注意到边界条件，如果边界条件是已知量（Dirichlet Boundary）则直接带入。离散化的
Laplace算子形如如下表达式：

\[
\nabla^2 f = \frac{f_{i+1,j} + f_{i-1,j} + f_{i,j+1} + f_{i,j-1} - 4f_{i,j}
}{h^2}
\]

对图片中所有的像素点使用Laplace算子，得到一个大型的线性方程组。使用Jacobi行列式进行迭代求解（Laplace Smoothing算法）
求解矩阵方程$\mathbf{Ax} = \mathbf{b}$,有直接法（LU分解），Jacobi迭代法，Gauss-Siebi迭代法。

\paragraph{Matting}

Matting就是一般所指的抠图。抠图使用的方法使用掩码和前景。已知Alpha算子对图像进行混合，需要区分前景和背景。混合方程（compositing equation）为：

\[
I_i = \alpha_i F_i + (1 - \alpha_i) B_i
\]

对RGB三通道分别成立：

\[
\begin{aligned}
I_i^{R} &= \alpha_i F_i^{R} + (1 - \alpha_i) B_i^{R} \\
I_i^{G} &= \alpha_i F_i^{G} + (1 - \alpha_i) B_i^{G} \\
I_i^{B} &= \alpha_i F_i^{B} + (1 - \alpha_i) B_i^{B}
\end{aligned}
\]

其中 $I_i$ 为观测像素值，$F_i$ 为前景色，$B_i$ 为背景色，$\alpha_i \in [0,1]$ 为Alpha掩码。每个像素只有3个方程却有 $F_i, B_i, \alpha_i$ 共7个未知数。
所以这个方程有无穷多个解。一种解决方法就是使用绿幕，令背景能够用已知的数值直接带入，减少未知数。这个方法的问题是，绿幕/蓝幕抠图由于漫反射可能会影响到前景。
在当代也可以用AI方法辅助解决。


\subsection{画图}

基本任务：连续空间图像进行光栅化转换，用离散模式进行表示，研究这些表示方法的高效性和有效性。

\subsubsection{直线的绘制}
首先以简单的算法为例：假设要绘制直线$y(x) = m x + b$。最直接的方法就是先建立坐标系，然后按照一定的步长遍历x轴，计算出y轴坐标之后进行取整。遍历x方向和y方向中函数增长较快的方向。因为k越小，浮点数的乘法越容易引发精度损失。

\begin{lstlisting}[style=MyCppStyle]
void line (int x0, int y0, int x1, int y1){
    float y = y0;
    float m = (y1 - y0)/ (float) (x1 - x0);
    int x;
    for(x=x0;x<=x1;x++) {
        float y= m*x + b;
        draw_pixel(x,Round(y));
    }
}
\end{lstlisting}

关注到上述Naive算法中出现了大量的浮点数运算，尤其是乘法，这是不经济的。利用上一步的数据，将乘积转换为加法迭代。关注到$y(x + 1) = y(x) + m$,将算法变为以下情况。

\begin{lstlisting}[style=MyCppStyle]
void line (int x0, int y0, int x1, int y1){
    float y = y0;
    float m = (y1 - y0)/ (float) (x1 - x0);
    int x;
    for(x=x0;x<=x1;x++) {
        draw_pixel(x,Round(y));
        y += m;
    }
}
\end{lstlisting}

把要求变得更加苛刻，可以直接避免浮点数运算，在累计过程中仅考虑整数，得到Brensenham绘制直线算法。这个算法的条件是，斜率范围处在$(0,1)$区间内。由于绘制的过程中，当前点如果是$(x, y)$,则下一个点要么横纵坐标都+1,要么只有横坐标+1。我们只需要判断直线的实际点离哪个更近，变成一个离散的二选一问题。

首先将直线的方程转换为隐函数$F(x，y)$,能够根据函数值的符号判断点和直线的位置关系，判断两个点的方法就是判断中点的位置，即看$F(x + 1, y + \frac{1}{2})$的符号。

\definecolor{figgreen}{RGB}{162,255,163}
\definecolor{figred}{RGB}{192,0,0}

\begin{center}
\begin{tikzpicture}[x=1cm,y=1cm]
  \fill[figgreen] (2.94,2.37) -- (9.23,4.77) -- (9.23,1.45) -- (2.94,1.45) -- cycle;
  \draw[blue,line width=0.06cm] (2.94,2.37) -- (9.22,4.73);
  \draw[line width=0.04cm,dash pattern=on 4.8pt off 3.7pt] (5.075,4.75) -- (5.075,1.46);
  \draw[line width=0.04cm,dash pattern=on 4.8pt off 3.7pt] (7.855,4.75) -- (7.855,1.46);
  \draw[line width=0.04cm,dash pattern=on 4.8pt off 3.7pt] (2.98,2.975) -- (9.11,2.975);
  \draw[line width=0.12cm] (3.75,1.83) -- (6.14,1.83);
  \fill (6.10,2.01) -- (6.48,1.83) -- (6.10,1.65) -- cycle;
  \draw[figred,line width=0.12cm] (3.75,1.83) -- (6.20,3.89);
  \fill[figred] (6.017,3.996) -- (6.468,4.113) -- (6.275,3.690) -- cycle;
  \draw[line width=0.035cm] (5.45,4.96) -- (6.20,4.45);
  \fill (6.240,4.568) -- (6.42,4.30) -- (6.105,4.370) -- cycle;
  \fill[figred] (3.75,1.83) circle (0.25);
  \fill (6.52,1.83) circle (0.24);
  \fill[figred] (6.50,4.14) circle (0.25);
  \draw[figred,line width=0.055cm,dash pattern=on 2.6pt off 2.6pt] (6.48,2.99) circle (0.235);
  \node at (4.07,0.97) {\Large$(x,y)$};
  \node at (6.92,0.97) {\Large$(x+1,y)$};
  \node at (7.40,4.83) {\Large$(x+1,y+1)$};
  \node at (3.60,5.27) {\Large\sffamily\bfseries Go here next?};
  \node[blue] at (1.71,2.12) {\large\sffamily\bfseries F(P)=0};
\end{tikzpicture}
\end{center}

要想避免浮点数得到隐函数，可以借用直线上已知点P,关注直线法向量和$\overrightarrow{PP_0}$的点乘结果，判断位置，从而构造出隐函数$F(P)$

\definecolor{figpink}{RGB}{255,197,207}

\begin{center}
\begin{tikzpicture}[x=1cm,y=1cm]
  \fill[figpink] (1.495,1.566) -- (1.495,4.075) -- (7.775,4.075) -- cycle;
  \fill[figgreen] (1.495,1.566) -- (7.775,4.075) -- (7.775,0.745) -- (1.495,0.745) -- cycle;
  \draw[line width=0.04cm] (1.495,0.745) rectangle (7.775,4.075);
  \draw[blue,line width=0.045cm] (2.23,1.90) -- (4.86,2.895);
  \draw[line width=0.03cm,dash pattern=on 0.6pt off 0.6pt] (1.495,1.566) -- (7.775,4.075);
  \fill[blue] (2.23,1.90) circle (0.125);
  \fill[blue] (4.86,2.895) circle (0.125);
  \draw[line width=0.04cm] (2.23,1.90) -- (3.00,3.265);
  \fill (3.00,3.265) circle (0.125);
  \draw[line width=0.04cm] (2.23,1.90) -- (2.55,1.33);
  \fill (2.468,1.383) -- (2.62,1.17) -- (2.574,1.327) -- cycle;
  \node at (2.645,3.47) {\Large$\mathbf{P}$};
  \node at (1.865,2.29) {\Large$\mathbf{P}_0$};
  \node at (5.34,2.83) {\Large$\mathbf{P}_1$};
  \node at (3.03,1.155) {\Large$\mathbf{N}$};
  \node at (5.75,3.695) {\LARGE$F<0$};
  \node at (6.77,2.92) {\LARGE$F>0$};
  \node[anchor=west,inner sep=0pt] at (1.36,5.265) {\fontsize{11.2}{13.2}\selectfont\sffamily implicit function: F(\textbf{P}) = \textbf{N}$\cdot$(\textbf{P} $-$ \textbf{P}\textsubscript{0})};
  \node[anchor=west,inner sep=0pt] at (1.37,4.715) {\fontsize{11.2}{13.2}\selectfont\sffamily\textbf{F} = 0 $\rightarrow$ \textbf{P} is on \textbf{L}};
\end{tikzpicture}
\end{center}

由于最后只关心和0的关系，所以可以不用归一化直线斜率等信息，直接进行整数运算。由于代入点的时候道理是一样的，可以把斜率乘以2,将向量的增量$(1, \frac{1}{2})$也变成整数进行计算。

最终得到：

\[
F(P) = 2N \cdot (P - P_0)
F(P + \Delta) = F(P) + 2N \cdot \Delta
\]

是一个可迭代的整数运算。从而可以进行循环，得到该运算的代码,借用了点乘来构造方程：

\begin{lstlisting}[style=MyCppStyle]
    int dx = 2*(x1-x0),
    int dy = 2*(y1-y0);
    int dydx = dy-dx;
    F = 0 + dy – dx/2;
    for (x=x0 ; x<=x1 ; x++) {
        draw_pixel(x, y);
        if (F<0) F += dy;
        else {y++; F += dydx;}
}
\end{lstlisting}

\subsubsection{曲线的绘制}

上述算法只要知道了曲线方程的隐函数表示，可以用类似比较大小的手段画曲线，但是由于二次甚至三次方程的求解较为困难，所以很难优化到整数级别的运算难度。

对于圆这样的图形，每一次迭代可以同时画出对称的8个点，实现类似于并行计算的效果。

\begin{lstlisting}[style=MyCppStyle]
    void draw_circle(int radius) 
    {
        int x = 0, y = radius;
        int d = 3-2*radius;
        while (y>x) {
            if (d<0) /* select East point next */
            d += 4*x + 6;
            else { /* select South-East point next */
            d += 4*(x-y) + 10;
            y--;}
        x++;
        draw_8_pts(x,y); /* draws point in each octant */
        }
}
\end{lstlisting}

\subsubsection{凸多边形的绘制}

凸多边形的判断并不容易，所以主要关注绘制三角形的算法。算法的基本原理是用一条直线自上而下进行扫描，一条一条地扫描，左右两侧的点由前述算法逐步确定。

对于一般的多边形可以进行三角化分割，归结到前面的算法。

具体来说，这套扫描转换流程需要依次完成以下几件事：先找出多边形的最高顶点与最低顶点，由此确定扫描线的起止范围；再把多边形的边沿着左右两侧分别组织成边表（左侧链与右侧链），按扫描线自上而下逐条处理；由于凸多边形每条扫描线只与边界相交两次，每条扫描线上只有一个待填充的区间（span），求出区间左右端点的横坐标$x_L, x_R$之后（端点也可以借助Bresenham算法顺带求出），把两者之间的像素全部填上颜色即可。相邻两条扫描线的端点之间满足

\[
x_L \mathrel{+}= \frac{dx_L}{dy_L}, \qquad x_R \mathrel{+}= \frac{dx_R}{dy_R}
\]

所以端点的推进同样是增量式的，与前面画直线的迭代完全一致。

上面这些环节中任何一处处理不当，都会在相邻多边形的公共边上留下痕迹：如果两侧链的取舍或者端点的取整规则不一致，公共边可能被填充两次，也可能一次都没有被填充，于是出现重叠（overlap）或者裂缝（crack）。

\subsubsection{凹多边形与三角化}

凹多边形的麻烦在于一条扫描线与它的交点可能多于两个，“每条扫描线只有一个区间”的前提不再成立。此时的做法是：求出该扫描线与多边形的全部交点，把它们按横坐标从左到右排序，再依次两两配对填充，这就是奇偶规则（parity rule）——排序后第1、2个交点之间是内部，第3、4个交点之间也是内部，奇数号交点为入点，偶数号交点为出点。也就是说，从某点向左引出的射线与多边形边界的交点个数为奇数时该点在内部，为偶数时在外部。

另一条思路是三角化（triangulation）：先把凹多边形分割成若干个三角形，再对每个三角形套用前面凸多边形的扫描转换流程。这样做的好处是内层循环永远是简单的单区间填充，不需要求交排序与奇偶判断；而只要共享同一条边的两个三角形对这条边的取整规则保持一致，拼接处既不会出现裂缝也不会重叠。代价是需要额外做一步三角化（对简单多边形可以用耳切法ear clipping，逐个切掉凸耳）。两种做法本质上是在“逐扫描线求交排序”与“事先分割、简化内层循环”之间做取舍。

\subsection{图像变换（Warping）}
（课堂上面只是做了科普，真正的内容远远少于下面的表述，表述是AI完成的）

在Image Warping中，进行图像的变形操作，比如对图片进行拉伸，扭曲，控制。具体而言就是整体或者局部的坐标系变换。

常见的方法是网格方法，将空间的网格进行拖动，挪移等操作。在PS软件中，实际上操作光标的中心是一个kernel，以相同的幅度和倍率进行平移和非线性形变。

对于RGB图像，一般会进行颜色插值，四边形和三角形的不同网格表示会导向不同的插值方法，非线性插值和线性插值都会出现在Warping中。

具体地，可以把图像看成定义在参数域上的函数：源图的一个正方形区域由四个角点参数化——左上$(u_0,v_0)$、右上$(u_1,v_0)$、左下$(u_0,v_1)$、右下$(u_1,v_1)$，Warping就是把这个参数正方形映射成目标图中的任意四边形，正方形的四个角落到四边形的四个角上，源图里的参数点$(u_i,v_j)$则落到四边形中对应的位置。图中的向量$\boldsymbol{u}$表示水平方向的映射，整张图可以理解为参数域到目标域的一次坐标变换，四边形的情形常取双线性映射，更一般（带透视）的情形则是单应（homography）变换。

实现上通常采用反向映射（inverse mapping）：遍历目标图像的每一个像素，用逆变换求出它在源图中对应的参数坐标$(u_i,v_j)$，再按这个坐标到源图里取样。之所以不反过来做正向映射（把源图的每个像素搬到目标位置），是因为正向映射下目标像素未必被覆盖，结果会出现空洞（holes），而多个源像素又可能落到同一个目标像素上造成重叠。

Warping的典型用途是几何矫正：把倾斜拍摄的画面（仰拍的埃菲尔铁塔，或者从侧面拍下来的蒙娜丽莎）重新映射回正视图，也就是所谓的keystone矫正；也可以在图像上覆盖一张控制网格，通过拖动网格的顶点做局部的非线性形变，这与PS里的变形/液化工具是同一个思路。

颜色插值的问题随之而来。把由两个三角形拼成的正方形（用红绿渐变示意）warp成三角形时，如果直接在屏幕空间对纹理坐标或者颜色做线性插值，得到的是“非透视校正”的结果：原本笔直的三角形公共边在结果图里会弯成一条曲线，两块三角形的拼接处发生错位，渐变也随之失真。原因是透视投影下屏幕上的等距并不对应纹理上的等距，参数坐标要先除以深度（齐次坐标里的$w$）才是线性的。正确的做法是透视校正插值（perspective-correct interpolation）：对$u/w$、$v/w$、$1/w$分别做线性插值，最后相除还原出$u,v$，这样三角形的公共边才会保持笔直。


\appendix

\section{数学知识补充}

\subsection{Fourier变换}

Fourier变换是图像处理中的常用数学规律，图像处理中涉及的主要是二维的Fourier变换，所以这里只探讨二维的情况，高维的情况可以自行推广。Fourier变换听起来很高级，但是实际上本质就是一个特殊的可逆线性变换，
(在图像处理中相当于变换坐标轴)，变换的目的类似于极坐标方程和直角坐标方程之间的转换。也就是对于同一个函数或者点，用不同的方式选取我们关注的信息进行表达。Fourier变换会有计算上和理解上的好处，我们马上就会看到。
\subsubsection{二维Fourier变换的定义}

\begin{definition}[二维Fourier变换离散版本]
    设定二维空间域中的函数为$f(x,y)$,则该函数的Fourier变换结果为
    \[
    F(u,v) = \sum_{x = 0}^{M - 1} \sum_{y = 0}^{N - 1} f(x,y) e^{-j2\pi\left(\frac{ux}{M} + \frac{vy}{N}\right)}
    \]
\end{definition}

连续版本的只需要把求和改写成二重积分，没有什么难度，而且和VCL处理关系不大，不予列举。现在我们需要明确DFT的含义。DFT变换之后得到的函数是一个复值函数，可以写成实部和虚部分离的形式：

\[
F(u,v) = R(u,v) + j I(u,v)
\]

二维Fourier变换的意义不太清楚，我们回忆一维Fourier变换的想法。对于一个信号函数，应当是由不同频率的正弦波叠加而成的。Fourier变换给出的关于频率$\omega$的函数说明了不同频率的波形情况。
将频率的概念推广到二维频域，则频域内的每一个点都对应了一个函数，说明了本频率空间内的信号情况。回看上式。$\sqrt{R^2 + I^2}$，被称之为模值，说明了频率的强弱也就是携带的能量（所谓高频低频指的就是这个）。
我们画的“频谱”也就是指“幅值”的频谱。而该复数的辐角$\theta = \arctan\frac{I}{R}$称之为相位。

\subsubsection{二维Fourier变换的性质}

Fourier变换不仅有分析上的好处，还有计算上的好处。首先二维Fouirer变换是一个线性变换，这点不难验证。从定义中可以看出，
这其实就是对二维矩阵进行了矩阵乘法。Fourier变换对应的矩阵是Fourier矩阵，定义如下：

\begin{definition}[Fourier矩阵]
    定义旋转因子：$$\omega_N = e^{-j\frac{2\pi}{N}}$$

    矩阵 $\mathbf{F}_N$ 的第 $k$ 行、第 $n$ 列（索引均从 $0$ 到 $N-1$）的元素为：
    
    $$\mathbf{F}_N[k, n] = \omega_N^{k \cdot n} = e^{-j\frac{2\pi}{N} kn}$$
\end{definition}

容易验证，这是一个酉矩阵，因此Fourier变换具有酉变换的所有性质。现在请我们回忆酉变换最重要的性质：

首先是可逆性，这意味着转换到频域之后不会有任何信息丢失。（也就是说对这个矩阵携带的信息进行了重新的分类而不改变内容）。

对于二维Fourier变换的逆变换有如下公式：

\begin{theorem}[Fourier变换的逆变换]
    设$f(x,y)$进行Fourier变换之后变成了$F(u,v)$,则逆变换形如以下表达式：
    \[
    f(x, y) = \frac{1}{MN} \sum_{u=0}^{M-1} \sum_{v=0}^{N-1} F(u,v) e^{j2\pi\left(\frac{ux}{M}+ \frac{vy}{N}\right)}
    \]
\end{theorem}

其次是保距性，由于是酉矩阵，变换前后点之间的距离是不会发生变化的。在研究频谱的视角下这意味着信号函数携带的总能量严格等于频域内各点所携带的总能量之和，此即著名的Parseval等式。能量
在VCL的信息处理中有极为重要的应用。

\begin{theorem}[Parseval等式]
    Parseval 等式写作：
    
    $$\sum_{x=0}^{M-1} \sum_{y=0}^{N-1} \vert{}f[x, y]\vert{}^2 = \frac{1}{MN} \sum_{u=0}^{M-1} \sum_{y=0}^{N-1} \vert{}F[u, v]\vert{}^2$$

    注：这时工程上常常采用的没有归一化的形式，严格的定义两面各自用一个常数归一化就可以。
\end{theorem}

这说明了每个频域点携带的能量，和原来每个像素点携带的能量是完全相等的。能量可以简单理解为信息量，这进一步说明Fourier变换前后信息没有丢失。

\subsubsection{Fourier变换的运算}
Fourier变换重要的运算性质有两个，首先是Fourier变换由于指数的性质容易降维拆分，在计算过程中可以直接先后计算两个一维Fourier变换而不需要计算二维的结构。正如以下公式：

\[
F(u, v) = \sum_{x=0}^{M-1} \left( \sum_{y=0}^{N-1} f(x, y) e^{-j 2\pi \frac{vy}{N}} \right) e^{-j 2\pi \frac{ux}{M}}
\]

首先对每行求，然后再对结果求，这样可以将计算复杂度由$O(M^2 N^2)$降到$O(MNlog(MN))$

其次是Fouier变换可以帮助计算卷积和高阶导数。卷积和高阶导数的计算是非常复杂的，然而经过Fourier变换，卷积可以退化为乘法计算，可以有效降低计算复杂度。具体关系如下，其中*表示做卷积。

\[
\mathcal{F} \{f * g\} = \mathcal{F} \{f\} \mathcal{F} \{g\} \qquad \mathcal{F} \{\frac{d^n f}{d x^n}\} = (j 2\pi u)^n \mathcal{F} \{u\}
\]

在进行复杂的卷积计算时，可以首先用DFT变换做乘法，然后变换回来，达到降低计算量的效果。

\section{Dithering算法补充}

Dithering算法的主要难点在于各种不同的噪声。我们应当知道噪声的原理，各种噪声的数学模型和生成方法。


\end{document}